% Blueprint: Character values of finite groups lie in cyclotomic fields
% Created by Numina

\begin{document}

\section{Statement of the main result}

Let $G$ be a finite group. The standard result we wish to formalize is
that the values of any finite-dimensional complex character of $G$ lie
in the cyclotomic field $\mathbb{Q}(\zeta_n)$, where
$n = \operatorname{exp}(G)$ is the exponent of $G$. Since
$\mathbb{Q}(\zeta_n)$ is a number field of characteristic zero, it
embeds (non-canonically) into $\mathbb{C}$, and we phrase the result
uniformly in terms of one such embedding.

\begin{theorem}[Character values lie in the cyclotomic field]
    \label{thm:brauer_character_in_cyclotomic}
    \uses{lem:trace_in_cyclotomic_range}
    Let $G$ be a finite group and let
    $n = \operatorname{exp}(G)$. There exists a ring homomorphism
    \[
        \varphi : \operatorname{CyclotomicField}(n,\mathbb{Q})
            \longrightarrow \mathbb{C}
    \]
    such that for every finite-dimensional complex representation
    $\rho : G \to \operatorname{GL}(V)$ and every $g \in G$,
    \[
        \operatorname{tr}\bigl(\rho(g)\bigr) \in \varphi(\operatorname{CyclotomicField}(n,\mathbb{Q})).
    \]
\end{theorem}
\begin{proof}
    \uses{lem:trace_in_cyclotomic_range, lem:cyclotomic_embeds_in_complex,
          lem:image_zeta_is_primitive_root, lem:exponent_pos}
    Let $n := \operatorname{exp}(G)$, which is a positive natural
    number by Lemma~\ref{lem:exponent_pos}. By
    Lemma~\ref{lem:cyclotomic_embeds_in_complex} there is a ring
    homomorphism
    $\varphi : \operatorname{CyclotomicField}(n,\mathbb{Q}) \to \mathbb{C}$.
    By Lemma~\ref{lem:image_zeta_is_primitive_root} the image
    $\zeta := \varphi\bigl(\operatorname{zeta} n\bigr)$ of the
    distinguished primitive $n$-th root of unity in
    $\operatorname{CyclotomicField}(n,\mathbb{Q})$ is a primitive
    $n$-th root of unity in $\mathbb{C}$. Fix such a $\varphi$ and
    $\zeta$. Now given $\rho$ and $g$, the order of $g$ in $G$
    divides $n$, so $\rho(g)^n = 1$ in
    $\operatorname{End}_{\mathbb{C}}(V)$. By
    Lemma~\ref{lem:trace_in_cyclotomic_range} applied to $\rho(g)$
    and the primitive $n$-th root $\zeta$, the trace
    $\operatorname{tr}(\rho(g))$ lies in $\varphi$'s range.
\end{proof}

\section{Auxiliary results about the cyclotomic field}

\begin{lemma}[Exponent of a finite group is positive]
    \label{lem:exponent_pos}
    For any finite group $G$, the exponent $\operatorname{exp}(G)$ is
    a positive natural number.
\end{lemma}
\begin{proof}
    Immediate from the Mathlib instance
    \texttt{Monoid.neZero\_exponent\_of\_finite}.
\end{proof}

\begin{lemma}[Cyclotomic field embeds in $\mathbb{C}$]
    \label{lem:cyclotomic_embeds_in_complex}
    \uses{lem:exponent_pos}
    For any positive natural number $n$, there exists a ring
    homomorphism
    \[
        \varphi : \operatorname{CyclotomicField}(n,\mathbb{Q})
            \longrightarrow \mathbb{C}.
    \]
\end{lemma}
\begin{proof}
    The field $\operatorname{CyclotomicField}(n,\mathbb{Q})$ is a
    number field (Mathlib instance
    \texttt{CyclotomicField.instNumberField}), hence algebraic over
    $\mathbb{Q}$ (\texttt{NumberField.isAlgebraic}). Since
    $\mathbb{C}$ is algebraically closed
    (\texttt{Complex.isAlgClosed}), Mathlib's
    \texttt{IsAlgClosed.lift} produces an algebra homomorphism
    $\operatorname{CyclotomicField}(n,\mathbb{Q}) \to_{\mathbb{Q}}
    \mathbb{C}$; the underlying ring homomorphism is the desired
    $\varphi$.
\end{proof}

\begin{lemma}[Image of $\zeta$ is a primitive $n$-th root of unity]
    \label{lem:image_zeta_is_primitive_root}
    \uses{lem:cyclotomic_embeds_in_complex, lem:exponent_pos}
    Let $n$ be a positive natural number and let
    $\zeta_0 = \operatorname{zeta}(n,\mathbb{Q},\operatorname{CyclotomicField}(n,\mathbb{Q}))$
    be the distinguished primitive $n$-th root of unity in
    $\operatorname{CyclotomicField}(n,\mathbb{Q})$.  For any
    injective ring homomorphism
    $\varphi : \operatorname{CyclotomicField}(n,\mathbb{Q}) \to \mathbb{C}$
    (such as the one of
    Lemma~\ref{lem:cyclotomic_embeds_in_complex}), the element
    $\varphi(\zeta_0) \in \mathbb{C}$ is a primitive $n$-th root of
    unity.
\end{lemma}
\begin{proof}
    By \texttt{IsCyclotomicExtension.zeta\_spec},
    $\zeta_0$ is a primitive $n$-th root of unity in
    $\operatorname{CyclotomicField}(n,\mathbb{Q})$. The lemma
    \texttt{IsPrimitiveRoot.map\_of\_injective} transports a
    primitive root along any injective ring (or monoid) homomorphism;
    a ring homomorphism between fields is automatically injective by
    \texttt{RingHom.injective}.
\end{proof}

\section{Trace of a finite-order endomorphism}

We isolate the linear-algebra core of the argument: the trace of a
finite-order $\mathbb{C}$-linear endomorphism on a finite-dimensional
$\mathbb{C}$-vector space is a sum of $n$-th roots of unity in
$\mathbb{C}$, and in particular lies in the subring generated by any
primitive $n$-th root of unity, which contains the image of any ring
homomorphism from $\operatorname{CyclotomicField}(n,\mathbb{Q})$ into
$\mathbb{C}$.

\begin{lemma}[$f^n = 1$ implies charpoly splits with roots in
    $\mu_n$]
    \label{lem:charpoly_roots_are_roots_of_unity}
    Let $V$ be a finite-dimensional $\mathbb{C}$-vector space, let
    $n \geq 1$, and let $f : V \to V$ be a $\mathbb{C}$-linear
    endomorphism with $f^n = 1$. Then the characteristic polynomial
    $\operatorname{charpoly}(f) \in \mathbb{C}[X]$ splits, and every
    multiset-element of $\operatorname{charpoly}(f).\operatorname{roots}$
    is an $n$-th root of unity in $\mathbb{C}$.
\end{lemma}
\begin{proof}
    \uses{lem:exponent_pos}
    Since $\mathbb{C}$ is algebraically closed
    (\texttt{Complex.isAlgClosed}), every polynomial over
    $\mathbb{C}$ splits (\texttt{IsAlgClosed.splits\_codomain}), so
    $\operatorname{charpoly}(f)$ splits.

    Let $\mu$ be any root of $\operatorname{charpoly}(f)$.  Then $\mu$
    is an eigenvalue of $f$
    (\texttt{Module.End.hasEigenvalue\_iff\_isRoot\_charpoly}). By
    \texttt{Module.End.HasEigenvalue.pow}, $\mu^n$ is an eigenvalue
    of $f^n$. But $f^n = 1$, so $\mu^n = 1$, i.e. $\mu$ is an
    $n$-th root of unity in $\mathbb{C}$ (\texttt{mem\_rootsOfUnity}).
\end{proof}

\begin{lemma}[Trace of a finite-order endomorphism is a sum of
    roots of unity]
    \label{lem:trace_eq_sum_roots_of_unity}
    \uses{lem:charpoly_roots_are_roots_of_unity}
    Let $V$ be a finite-dimensional $\mathbb{C}$-vector space, let
    $n \geq 1$, and let $f : V \to V$ be a $\mathbb{C}$-linear
    endomorphism with $f^n = 1$. Then there exists a multiset
    $\mu_1, \ldots, \mu_k$ of $n$-th roots of unity in $\mathbb{C}$
    such that
    \[
        \operatorname{tr}(f) = \sum_{i=1}^{k} \mu_i.
    \]
\end{lemma}
\begin{proof}
    \uses{lem:charpoly_roots_are_roots_of_unity}
    By
    \texttt{Module.End.trace\_eq\_sum\_roots\_charpoly\_of\_splits}
    and the fact that $\operatorname{charpoly}(f)$ splits
    (Lemma~\ref{lem:charpoly_roots_are_roots_of_unity}),
    \[
        \operatorname{tr}(f) =
            \operatorname{charpoly}(f).\operatorname{roots}.\operatorname{sum}.
    \]
    By Lemma~\ref{lem:charpoly_roots_are_roots_of_unity} every root
    appearing in this multiset is an $n$-th root of unity, so the
    multiset itself serves as the required $\mu_1, \ldots, \mu_k$.
\end{proof}

\begin{lemma}[$n$-th roots of unity lie in the subring generated by
    a primitive root]
    \label{lem:nth_root_in_powers}
    \uses{lem:image_zeta_is_primitive_root}
    Let $\zeta \in \mathbb{C}$ be a primitive $n$-th root of unity
    and let $\mu \in \mathbb{C}$ satisfy $\mu^n = 1$. Then there
    exists $k \in \mathbb{N}$ with $\mu = \zeta^k$.
\end{lemma}
\begin{proof}
    The element $\mu$ lies in $\operatorname{rootsOfUnity}(n)$
    by \texttt{mem\_rootsOfUnity}, and Mathlib's
    \texttt{IsPrimitiveRoot.eq\_pow\_of\_mem\_rootsOfUnity} (or the
    equivalent \texttt{IsPrimitiveRoot.eq\_pow\_of\_pow\_eq\_one})
    yields the required exponent $k$.
\end{proof}

\begin{lemma}[Trace lies in the range of $\varphi$]
    \label{lem:trace_in_cyclotomic_range}
    \uses{lem:trace_eq_sum_roots_of_unity, lem:image_zeta_is_primitive_root,
          lem:nth_root_in_powers, lem:cyclotomic_embeds_in_complex,
          lem:exponent_pos}
    Let $n \geq 1$, let
    $\varphi : \operatorname{CyclotomicField}(n,\mathbb{Q}) \to
    \mathbb{C}$ be a ring homomorphism, and let
    $\zeta := \varphi(\operatorname{zeta}\,n)$, which is a primitive
    $n$-th root of unity in $\mathbb{C}$
    (Lemma~\ref{lem:image_zeta_is_primitive_root}).  For any
    finite-dimensional $\mathbb{C}$-vector space $V$ and any
    $\mathbb{C}$-linear endomorphism $f : V \to V$ with $f^n = 1$,
    \[
        \operatorname{tr}(f) \in \varphi(\operatorname{CyclotomicField}(n,\mathbb{Q})).
    \]
\end{lemma}
\begin{proof}
    \uses{lem:trace_eq_sum_roots_of_unity, lem:nth_root_in_powers,
          lem:image_zeta_is_primitive_root}
    By Lemma~\ref{lem:trace_eq_sum_roots_of_unity},
    $\operatorname{tr}(f) = \sum_{i} \mu_i$ for some $n$-th roots of
    unity $\mu_i \in \mathbb{C}$. By
    Lemma~\ref{lem:nth_root_in_powers} each $\mu_i$ equals
    $\zeta^{k_i}$ for some $k_i \in \mathbb{N}$. Since
    $\zeta = \varphi(\operatorname{zeta}\,n)$ lies in
    $\varphi$'s range and $\varphi$'s range is a subring
    (\texttt{RingHom.range}, closed under sums and powers), so does
    each $\mu_i$, and hence so does $\operatorname{tr}(f)$.
\end{proof}

\section{Bridging back to representations}

We finally specialize the previous linear-algebra statement to a
representation $\rho$ of a finite group $G$.

\begin{lemma}[Finite group element raised to the exponent is the
    identity in any representation]
    \label{lem:rho_g_pow_exp}
    \uses{lem:exponent_pos}
    Let $G$ be a finite group, let $V$ be a complex vector space,
    and let $\rho : \operatorname{Representation}\,\mathbb{C}\,G\,V$.
    Let $n := \operatorname{exp}(G)$. Then for every $g \in G$,
    $(\rho\,g)^n = 1$ in $\operatorname{End}_{\mathbb{C}}(V)$.
\end{lemma}
\begin{proof}
    By \texttt{Monoid.pow\_exponent\_eq\_one}, $g^n = 1$ in $G$.
    Since $\rho$ is a monoid homomorphism $G \to
    \operatorname{End}_{\mathbb{C}}(V)$ (\texttt{Representation} is
    defined as a \texttt{MonoidHom} to the endomorphism monoid),
    \texttt{MonoidHom.map\_pow} and \texttt{MonoidHom.map\_one} give
    $(\rho\,g)^n = \rho(g^n) = \rho 1 = 1$.
\end{proof}

The proof of Theorem~\ref{thm:brauer_character_in_cyclotomic} above
then assembles
Lemma~\ref{lem:rho_g_pow_exp} and
Lemma~\ref{lem:trace_in_cyclotomic_range} (with the
$\varphi$ produced by
Lemma~\ref{lem:cyclotomic_embeds_in_complex}).

\end{document}
